% ============================================================================
% Calibration
% ============================================================================

Accurate ranging requires removing two classes of bias:
a bulk per-board constant offset (antenna delay) and a
position-dependent residual (multipath geometry).
We address these in two independent tiers, and document a third calibration
failure mode that motivated this design.

\subsection{Tier 1: Hardware Antenna Delay Calibration}

Each DW1000 has an internal timestamp latency between first-path arrival and
the hardware capture event.
This latency, stored as a 16-bit register in DW1000 ticks
($1\,\text{tick} \approx \SI{15.65}{\pico\second} \approx \SI{4.69}{\milli\metre}$),
is board-specific.
The factory default of 16384 ticks leaves uncalibrated boards reading
approximately 30--50\,cm too long.

\textbf{Calibration procedure.}
The host sends a \texttt{HWCALIB,<id>,<true\_dist\_mm>} UDP command.
The tag executes a binary search: in each of $I = 14$ iterations it pushes
a candidate delay $d_c$ to the target anchor via a \texttt{MSG\_ANT\_DELAY}
UWB relay, collects $S = 50$ range samples, and bisects based on whether the
mean range exceeds the known true distance.
After $I$ iterations the resolution is $< 0.07\,\text{ticks} \approx 0.3$\,mm.
The final delay is pushed to the anchor's NVS via \texttt{Preferences}, so
it survives reflashing.

\textbf{Calibrated values.}
Table~\ref{tab:ant_delays} lists the calibrated antenna delays for all nodes,
obtained at a reference distance of \SI{3.0}{\metre}.

\begin{table}[h]
  \centering
  \caption{Calibrated antenna delays (DW1000 ticks $\approx$ \SI{4.69}{\milli\metre}/tick).
  Reference distance: \SI{3.0}{\metre}. PRF: 64\,MHz, 110\,kbps.}
  \label{tab:ant_delays}
  \begin{tabular}{lccc}
    \toprule
    Node & Address & Delay (ticks) & $\Delta$ from factory (ticks) \\
    \midrule
    Anchor 1 & \texttt{0x01} & 16422 & $+38$ ($\approx\SI{178}{\milli\metre}$) \\
    Anchor 2 & \texttt{0x02} & 16407 & $+23$ ($\approx\SI{108}{\milli\metre}$) \\
    Anchor 3 & \texttt{0x03} & 16394 & $+10$ ($\approx\SI{47}{\milli\metre}$)  \\
    Anchor 4 & \texttt{0x04} & 16522 & $+138$ ($\approx\SI{647}{\milli\metre}$) \\
    Anchor 5 & \texttt{0x05} & 16410 & $+26$ ($\approx\SI{122}{\milli\metre}$) \\
    Tag 1    & \texttt{0xF0} & 16360 & $-24$ ($\approx\SI{-113}{\milli\metre}$) \\
    Tag 2    & \texttt{0xF1} & 16473 & $+89$ ($\approx\SI{418}{\milli\metre}$) \\
    \bottomrule
  \end{tabular}
\end{table}

The calibrated delays span 128\,ticks across anchors ($+10$ to $+138$ above
factory default 16384), confirming that per-board calibration is essential;
a shared default would introduce systematic ranging errors of up to
\SI{647}{\milli\metre} (\SI{64.7}{\centi\metre}) on Anchor~4.

\subsection{Tier 2: Spatial Error Map}
\label{sec:tier2}

After Tier-1 calibration, residual ranging errors remain due to
position-dependent multipath: the channel impulse response varies with
tag location because different wall-reflection paths are constructive or
destructive at different points in the room.
Offline analysis of calibration logs (see \S\ref{sec:failure_mode}) showed
this residual is non-monotonic in range and cannot be modelled by a
per-anchor polynomial correction.
A \emph{spatial map}---a correction $c_i(\mathbf{p})$ per anchor as a
function of tag position---matches the physics.

\textbf{Grid capture.}
The tag is placed at each of $G$ grid points (default: $G = 9$, a
$3 \times 3$ grid at $\pm\SI{1.5}{\metre}$ around the anchor centroid).
At each point $\mathbf{g}_k$, $S_k$ raw distance samples are collected for
each visible anchor~$i$; samples with gap $> \SI{6}{\decibel}$ are
rejected (NLOS capture gate).
The residual at grid point~$k$ for anchor~$i$ is
\begin{equation}
  r_{i,k} = \bar{\rho}_{i,k} - \|\mathbf{g}_k - \mathbf{a}_i\|_2,
  \label{eq:errormap_residual}
\end{equation}
where $\bar{\rho}_{i,k}$ is the mean over accepted samples.

\textbf{Runtime correction.}
Given the rover's previous-sweep EKF position $\hat{\mathbf{p}}_{t-1}$
(one-step lag of $\sim\SI{100}{\milli\second}$, negligible at rover speed),
the correction for anchor~$i$ is the inverse-distance-weighted (IDW) mean
over the four nearest grid points:
\begin{equation}
  c_i(\hat{\mathbf{p}}_{t-1}) =
  \frac{\displaystyle\sum_{k \in \mathcal{N}(4)} w_k\,r_{i,k}}
       {\displaystyle\sum_{k \in \mathcal{N}(4)} w_k},
  \quad w_k = \frac{1}{d_k^2 + \varepsilon^2},
  \label{eq:idw}
\end{equation}
$d_k = \|\hat{\mathbf{p}}_{t-1} - \mathbf{g}_k\|$, $\varepsilon = \SI{0.05}{\metre}$.
The correction is linearly faded to zero beyond $1.5 \times s_{\text{grid}}$
from the nearest grid point and clamped to $\pm\SI{0.30}{\metre}$.
The corrected range is $\rho_i^{\text{corr}} = \rho_i - c_i$.

\subsection{Documented Calibration Failure Mode}
\label{sec:failure_mode}

A prior version of the system (version~v2/v3 calibration JSON) applied a
Level-2 software bias \emph{simultaneously} with the Level-1 hardware delay
push.
Concretely: the auto-calibration routine computed a per-anchor bias $b_0$
from reference-position ranging, then (a)~derived $\Delta = b_0/4.69$
DW1000 ticks and pushed this to the anchor NVS, and (b)~subtracted $b_0$
from every subsequent range measurement in the host pipeline.
After the hardware push removed $b_0$ from the physics, the software
subtracted $b_0$ again, producing an over-correction of $2 b_0$.
With uncalibrated boards showing $b_0 \in [+445, +708]$\,mm, the
double-correction inverted the bias---the system behaved worse than no
calibration at all.

A second bug was also present: the delay push used the sign convention
\texttt{new\_delay = current\_delay $-\Delta$}, whereas the correct sign is
\texttt{+$\Delta$} (more antenna delay shortens the reported distance;
removing a positive bias requires adding delay).
Rather than patching these bugs, we removed the Level-2 software-bias path
entirely and replaced it with the position-dependent Tier-2 spatial map,
which has no coupling to the Tier-1 hardware correction.
This is described in full in a supplementary technical report
\todoref{cite our calibration redesign document (docs/calibration\_redesign\_2tier.md)
as supplementary material or arXiv technical note before submission}.
The v4 \texttt{calibration.json} schema retains only Tier-1 traceability;
old v2/v3 files load without error but their Level-2 fields are discarded
with a warning.
